% =============================================================================================
\chapter{Introduction}\label{chap:intro}

\section{Contexte}

Une banque reçoit chaque année des demandes de crédit d'investissement émanant de petites et moyennes entreprises :
un garage qui veut s'équiper, une clinique qui s'agrandit, un transporteur qui renouvelle sa flotte. La banque est le
\emph{financeur} ; les projets appartiennent à des entités extérieures. Ses ressources sont limitées : un montant
global est alloué sur une période donnée, chaque projet doit respecter des montants minimum et maximum, et la banque
souhaite financer des projets variés plutôt que concentrer son risque sur un seul secteur.

Pour chaque dossier, plusieurs dimensions comptent et se contredisent souvent. Un projet très rentable peut être
risqué ; un projet sûr peut créer peu de valeur ou peu d'emplois ; un projet bien garanti peut être mal noté sur le
plan environnemental. La décision est donc par nature \emph{multicritère}. Elle est aussi
\emph{réglementée} : la banque doit pouvoir expliquer un refus et documenter son processus d'octroi, comme le
rappellent les orientations de l'Autorité bancaire européenne sur l'octroi et le suivi des prêts~\cite{eba2020}.

Le sujet proposé est de s'appuyer sur l'\emph{IA agentique} : des agents spécialisés analysent les dossiers,
évaluent chaque critère et proposent un classement argumenté des projets, la décision finale restant aux décideurs
de la banque. Le choix imposé est un modèle de classement en trois catégories : \og à financer \fg{},
\og à financer si\dots{} \fg{} (sous conditions) et \og rejeté \fg{}.

\section{Problématique}

\begin{encadre}[Problématique]
Comment combiner une estimation du risque de défaut apprise sur des données réelles, des critères financiers et
extra-financiers et des contraintes de budget pour proposer, de façon transparente et vérifiable, un classement et
une catégorie pour chaque projet, tout en utilisant des modèles de langage sans leur confier ni calcul ni décision ?
\end{encadre}

Trois difficultés structurent cette question. La première est \emph{statistique} : le défaut est rare (environ 4\,\%
des prêts récents), les données mélangent des périodes de crise et de calme, et une variable peut contenir l'issue
du prêt sans qu'on le voie (fuite de données). La deuxième est \emph{décisionnelle} : les critères sont hétérogènes
(une probabilité, des montants, des ratios, un score) et une mauvaise note sur un critère essentiel ne doit pas être
rachetée par de bonnes notes ailleurs. La troisième est \emph{de confiance} : un modèle de langage rédige bien mais
peut inventer des chiffres, ce qui est inacceptable dans un dossier de crédit.

\section{Objectifs}

\begin{enumerate}
  \item Construire une variable de défaut à cinq ans fiable à partir de prêts réels et ouverts (programme SBA 7(a)),
        sans fuite d'information.
  \item Comparer un modèle de référence simple et des modèles plus avancés, les calibrer, mesurer leurs
        performances sur des années futures et expliquer leurs prédictions.
  \item Évaluer un portefeuille de projets candidats sur huit critères par une méthode multicritère non
        compensatoire (ELECTRE III et ELECTRE TRI), vérifiée par une seconde méthode (PROMETHEE II), avec des poids
        cohérents (AHP) et un filtre éliminatoire.
  \item Sélectionner les projets à financer sous un budget global, des bornes de montant et un plafond sectoriel.
  \item Mesurer la robustesse des décisions aux paramètres de la méthode et aux données simulées.
  \item Ajouter des agents fondés sur un modèle de langage qui rédigent et expliquent sans jamais calculer.
  \item Livrer une solution déployée, réutilisable sur d'autres jeux de projets, reproductible et documentée.
\end{enumerate}

\section{Principe directeur}

Le projet repose sur une séparation stricte des rôles, retenue dès la conception :

\begin{encadre}[Qui fait quoi]
\textbf{L'apprentissage automatique ne calcule que le risque.} \textbf{La méthode multicritère classe et catégorise.}
\textbf{L'agent d'arbitrage respecte le budget.} \textbf{Le modèle de langage rédige et répond, il ne calcule jamais.}
\textbf{Le comité d'investissement fixe les paramètres en amont et décide en aval.}
\end{encadre}

Cette séparation a deux avantages. Chaque composant peut être testé et justifié séparément, et chaque chiffre d'une
recommandation peut être retracé jusqu'aux données et aux paramètres qui l'ont produit.

\section{Contributions}

\begin{itemize}
  \item Un pipeline reproductible sur 910\,033 prêts SBA, avec un audit des données, une séparation temporelle et la
        détection d'une fuite de données non documentée (la durée exacte du prêt).
  \item Une comparaison de trois modèles qui chiffre le \emph{prix de l'interprétabilité} : 0,022 point d'AUROC entre
        le meilleur modèle interprétable (EBM) et un modèle de boosting opaque (LightGBM).
  \item Un moteur multicritère implémenté et testé à la main (AHP, filtre, ELECTRE III avec distillation, ELECTRE TRI,
        PROMETHEE II, sensibilité sur 200 scénarios et 20 tirages de simulation).
  \item Un agent d'arbitrage sous budget, avec deux règles de sélection comparées.
  \item Deux agents fondés sur un modèle de langage, dont un conseiller qui appelle dix outils de calcul, avec un
        vérificateur automatique des nombres et des noms de méthodes.
  \item Une application Streamlit déployée et un moteur générique qui évalue n'importe quel jeu de projets.
\end{itemize}

\section{Plan du rapport}

Le chapitre~\ref{chap:theorie} présente les fondements théoriques de toutes les méthodes utilisées. Le
chapitre~\ref{chap:donnees} décrit les données, leur nettoyage et l'analyse exploratoire. Le
chapitre~\ref{chap:methodo} expose l'architecture du système et la méthodologie. Le chapitre~\ref{chap:resultats}
présente les résultats. Le chapitre~\ref{chap:solution} décrit la solution déployée. Le
chapitre~\ref{chap:discussion} discute les choix, les limites et les perspectives, avant la conclusion.
